% ---------------------------------------------------------------------
%  Appendix C : An Annotated File Index of the Source Tree
% ---------------------------------------------------------------------
\chapter{An Annotated File Index of the Source Tree}\label{app:file-index}

The preceding chapters walked the pipeline \emph{by stage} --- the
order in which data flows from a model description to a plotted cumulant.
This appendix walks the same code \emph{by directory}. It is a map, not a
tutorial: one line per module, what that module is responsible for, and
which chapter of this manual opens it up. If you arrive at the repository
cold, knowing the physics but not the layout, read this appendix first to
find the file you want, then jump to the chapter the table points you at.

The repository has a deliberate two-layer split, and understanding it is
the single most useful orientation you can have:

\begin{description}[style=nextline,leftmargin=1.4em]
  \item[\file{msrjd/} --- the engine.] The model-agnostic physics core. It
        knows about \msrjd{} actions, Feynman diagrams, propagators, loop
        integrals, and symmetry factors, but it knows \emph{nothing} about
        any particular model. It is pure \code{import}-able Python (running
        under SageMath) with no notebook dependencies.
  \item[\file{pipeline/} --- the driver.] The orchestration layer that
        sequences the engine's stages, threads data dictionaries between
        them, caches the expensive ones to disk, and exposes a single
        master function, \code{compute\_cumulants}. It also hosts the
        \emph{authoring} surface: the \code{TheoryBuilder} API, the
        text-to-Sage compiler, the colored-noise rewriter, and the
        \code{ipywidgets} theory-builder UI.
  \item[\file{models/} --- the validators.] Independent direct simulators
        (Euler--Maruyama SDE integrators, Poisson--Hawkes steppers,
        spectral PDE solvers) and the cumulant estimators that turn their
        output into the same correlator the engine predicts. These share
        \emph{no code} with the engine; they exist so every prediction can
        be checked against an experiment.
  \item[\file{theories/} --- the models.] The declarative
        \file{*.theory.py} files: one physical model each, the
        version-controllable single source of truth.
  \item[\file{notebooks/} --- the front desk.] The user-facing convenience
        layer \code{daedalus.py} and the example, template, and validation
        notebooks that drive it.
\end{description}

\begin{note}
\textbf{How to read the tables.} Each table lists files under one
directory. The first column is the file (always relative to the repository
root, the navy monospace convention of this manual). The second column is
its role in one phrase. The third column, ``Ch.'', is the chapter that
covers it in depth --- a dash means the file is peripheral (a stub, an
oracle-only cross-check, or a reference asset) and is documented here only.
The roles below are taken from each module's own docstring and verified
against its code; where a module's docstring and its behaviour disagree,
the table reflects the \emph{behaviour}, and the discrepancy is flagged in
Appendix~D (Known Issues and Open Questions).
\end{note}

\begin{gotcha}
\textbf{Two things named ``serialize'' and two things named
``propagator''.} The tree contains deliberate near-collisions that trip up
newcomers. There are \emph{two} serializers ---
\file{pipeline/theory\_serialize.py} (the \file{.theory.py} \emph{text}
format) and \file{msrjd/core/serialize.py} (the \code{.sobj} \emph{binary}
cache) --- and \emph{two} files called \code{propagator.py}, of which only
\file{pipeline/\_propagator.py} is live; \file{msrjd/core/propagator.py} is
a docstring-only stub. The tables call out each one. When you grep for a
name, check which of the pair you have landed on.
\end{gotcha}

% =====================================================================
\section{The engine: \texttt{msrjd/}}
% =====================================================================

\file{msrjd/} is the physics core, organised into five sub-packages:
\file{core/} (the symbolic action, vertices, propagator-time-domain, and
caching primitives), \file{enumeration/} (diagram topology generation),
\file{diagrams/} (typing, causality, symmetry), \file{integration/} (the
two loop-integral back ends, temporal and spatial), and a top-level
\file{fork\_safety.py} guard shared across the parallel paths.

\subsection{\texttt{msrjd/core/} --- the symbolic core}

This is where a model dict becomes a symbolic \msrjd{} action, gets
Taylor-expanded around the saddle, and is sorted into the free part (the
inverse propagator) and the interaction vertices. Everything here is a
SageMath symbolic-ring object; nothing is bound to a number yet.

\begin{center}
\begin{longtable}{@{}>{\ttfamily\footnotesize\raggedright}p{0.31\textwidth} >{\raggedright}p{0.50\textwidth} >{\centering}p{0.05\textwidth}@{}}
\toprule
\normalfont\textbf{file} & \textbf{role} & \textbf{Ch.} \tabularnewline
\midrule
\endhead
\bottomrule
\endfoot
msrjd/core/field\_theory.py
  & The \code{FieldTheory} expander: builds the symbolic namespace,
    evaluates the action lambda, resolves \code{Conv(...)} atoms, injects
    correlated-noise cumulants, runs the multivariate Taylor expansion, and
    bigrade-classifies the result into \code{ft.\_by\_tp} sectors. The
    front end of the whole pipeline.
  & \ref{ch:fieldtheory-core} \tabularnewline
\addlinespace
msrjd/core/vertices.py
  & Decomposes each bigrade sector polynomial into individual typed
    monomials: \code{VertexType} (interaction), \code{SourceType} (noise),
    and the kernel-carrying subclasses \code{ConvVertexType} (conductance
    $g(\tau)$ kernel) and \code{NoiseSourceType} (non-local $\kappa(\tau)$
    kernel).
  & \ref{ch:fieldtheory-core} \tabularnewline
\addlinespace
msrjd/core/convolution.py
  & The formal \code{Conv(kernel, field)} operator and its reduction rules
    --- how a registered synaptic kernel multiplied by a field is held
    symbolically until the propagator step can resolve it.
  & \ref{ch:fieldtheory-core} \tabularnewline
\addlinespace
msrjd/core/serialize.py
  & Save/load of an \emph{expanded} theory to a directory
    (\code{metadata.json} $+$ \code{symbolic\_data.sobj}). The
    \code{.sobj}-cache serializer --- not the \file{.theory.py} text format
    (that is \file{pipeline/theory\_serialize.py}).
  & \ref{ch:fieldtheory-core},\,\ref{ch:caching} \tabularnewline
\addlinespace
msrjd/core/cache.py
  & \code{PipelineCache}: the low-level stage-keyed \code{.sobj} store both
    higher-level caches build on. Keys each pipeline stage (enumeration,
    filtering, typing, dedup, grouping) by $(k,\ell)$ so a notebook restart
    skips the expensive recomputation.
  & \ref{ch:caching} \tabularnewline
\addlinespace
msrjd/core/propagator.py
  & \textbf{Stub only.} A ten-line docstring marked ``Status: Not yet
    extracted --- logic currently lives in the demo notebook''. Imported by
    nothing on the working path; the live propagator code is
    \file{pipeline/\_propagator.py}.
  & --- \tabularnewline
\end{longtable}
\end{center}

\subsection{\texttt{msrjd/enumeration/} --- diagram topology}

Given only the two integers $k$ (correlator order) and $\ell$ (loop order),
this package produces the complete, deduplicated list of bare graph
topologies --- \term{prediagrams} --- that could contribute, before any
physics is attached. It is theory-agnostic: it never looks at the action.

\begin{center}
\begin{longtable}{@{}>{\ttfamily\footnotesize\raggedright}p{0.31\textwidth} >{\raggedright}p{0.50\textwidth} >{\centering}p{0.05\textwidth}@{}}
\toprule
\normalfont\textbf{file} & \textbf{role} & \textbf{Ch.} \tabularnewline
\midrule
\endhead
\bottomrule
\endfoot
msrjd/enumeration/\newline loop\_diagram\_enumeration.py
  & The three-stage topology generator: nonisomorphic trees $\to$
    topologies (add $\ell$ edges) $\to$ prediagrams (orient every edge),
    each stage a filter-plus-deduplicate. Calls the \code{nauty} graph
    canonicaliser through Sage for the isomorphism rejection.
  & \ref{ch:enumeration} \tabularnewline
\addlinespace
msrjd/enumeration/degree\_scan.py
  & Build Phase C, the feedback bridge. Scans the enumerated prediagrams
    for their maximum vertex degree and, if a diagram needs a higher-degree
    vertex than the action was expanded to, triggers a re-expansion via
    \code{available\_degrees}.
  & \ref{ch:enumeration} \tabularnewline
\end{longtable}
\end{center}

\subsection{\texttt{msrjd/diagrams/} --- typing, causality, symmetry}

This package turns a bare oriented graph into a fully labelled Feynman
diagram you can integrate: it decorates edges with propagators and vertices
with monomials (typing), prunes the acausal ones, deduplicates, and
attaches the symmetry factor $\Sym(\Gamma)$.

\begin{center}
\begin{longtable}{@{}>{\ttfamily\footnotesize\raggedright}p{0.31\textwidth} >{\raggedright}p{0.50\textwidth} >{\centering}p{0.05\textwidth}@{}}
\toprule
\normalfont\textbf{file} & \textbf{role} & \textbf{Ch.} \tabularnewline
\midrule
\endhead
\bottomrule
\endfoot
msrjd/diagrams/filter.py
  & Build Phase D. A cheap set-membership check that discards prediagrams
    whose vertex-degree signature no vertex or source in the theory could
    supply --- run \emph{before} the expensive typing step.
  & \ref{ch:type-assignment} \tabularnewline
\addlinespace
msrjd/diagrams/\newline type\_assignment.py
  & Build Phase E, the heart of the subsystem. A backtracking
    constraint-satisfaction engine that emits every valid
    \code{TypedDiagram}: which monomial sits at each vertex, which external
    field on each leg, and which propagator $G_{ij}$ rides each directed
    edge.
  & \ref{ch:type-assignment} \tabularnewline
\addlinespace
msrjd/diagrams/causality.py
  & Build Phase F. Prunes typed diagrams that violate the retarded boundary
    condition --- structurally, the graph must be a DAG (no directed cycle);
    analytically, all propagator poles in the upper half $\omega$-plane.
  & \ref{ch:type-assignment} \tabularnewline
\addlinespace
msrjd/diagrams/symmetry.py
  & Build Phase G. Deduplicates typed diagrams via a complete
    graph-isomorphism invariant (Sage canonical labelling) and computes the
    symmetry factor
    $\Sym(\Gamma)=\big(\prod_v\prod_\ell n_{v,\ell}!\big)/|\Aut_{\text{fixed-ext}}(\Gamma)|$.
    Also hosts \code{external\_wick\_compensation} and
    \code{vertex\_role\_signature}, used by the integrators.
  & \ref{ch:type-assignment},\,\ref{ch:grouped-phasej} \tabularnewline
\addlinespace
msrjd/fork\_safety.py
  & The canonical fork-safety guard shared by every fork-based
    multiprocessing path. Detects ``am I inside a macOS Jupyter/ZMQ
    kernel?'' and, if so, silently degrades a parallel path to serial to
    avoid the fork-after-Cocoa/BLAS kernel crash.
  & \ref{ch:type-assignment} \tabularnewline
\end{longtable}
\end{center}

\subsection{\texttt{msrjd/integration/time\_domain/} --- the temporal loop integrator}

This is ``Phase J'': it takes each typed diagram's product of retarded
propagators and integrates over the internal vertex times. Because every
retarded propagator is a sum of decaying exponentials, the integral over a
causal-ordering polytope has a closed form, so the engine almost never
falls back to brute quadrature.

\begin{center}
\begin{longtable}{@{}>{\ttfamily\footnotesize\raggedright}p{0.31\textwidth} >{\raggedright}p{0.50\textwidth} >{\centering}p{0.05\textwidth}@{}}
\toprule
\normalfont\textbf{file} & \textbf{role} & \textbf{Ch.} \tabularnewline
\midrule
\endhead
\bottomrule
\endfoot
msrjd/integration/\newline time\_domain/final\_integral.py
  & The 5{,}269-line core. \code{integrate\_diagram} and the family of
    closed-form mode-sum evaluators (\code{\_integrate\_1d\_polytope\_modesum},
    \code{\_integrate\_2d\_polygon\_modesum},
    \code{\_integrate\_nd\_polytope\_poset\_modesum}), the \code{EdgeModeSum}
    representation, and the pole-tuple enumerator. The numerical heart of
    the temporal pipeline.
  & \ref{ch:phasej-temporal} \tabularnewline
\addlinespace
msrjd/integration/\newline time\_domain/pipeline.py
  & The orchestration/batching layer over the per-diagram core:
    \code{compute\_correction\_td} loops the integrator once per typed
    diagram and assembles the per-loop-order correction. Carries the
    fork-guarded \code{total\_C\_batch} parallel path.
  & \ref{ch:phasej-temporal} \tabularnewline
\addlinespace
msrjd/integration/\newline time\_domain/propagator\_td.py
  & Turns frequency-domain propagator data (poles, residues,
    \code{D\_delta}) into the time-domain matrix $G(t)$ = pole-residue sum
    $+\;\delta(t)$ part, with per-edge lookups \code{G\_t\_entry} and
    \code{G\_t\_delta\_coeff} the integrators call.
  & \ref{ch:propagator},\,\ref{ch:phasej-temporal} \tabularnewline
\addlinespace
msrjd/integration/\newline time\_domain/grouped\_integral.py
  & The grouped evaluator \code{integrate\_grouped\_diagram} and its
    closed-form helpers: evaluates all typed diagrams sharing one prediagram
    together, summing the per-edge propagator-index choices into one
    mode-sum instead of re-integrating each.
  & \ref{ch:grouped-phasej} \tabularnewline
\addlinespace
msrjd/integration/\newline time\_domain/subgraph.py
  & \textbf{Near-stub.} \code{identify\_loop\_subgraphs} returns \code{[]}
    (a never-finished loop-reduction route). Still re-exported by the
    package \code{\_\_init\_\_} and touched by a test, so it is a live import
    surface with a dead body.
  & --- \tabularnewline
\end{longtable}
\end{center}

\subsection{\texttt{msrjd/integration/spatial/} --- the spatial loop integrator}

For a stochastic \emph{partial} differential equation the engine does the
momentum integral first and analytically: every diagram, any $k$, any
$\ell$, any dimension $d$, reduces to the \emph{same} Symanzik integral and
is evaluated by the \emph{same} code. Several files in this package are
banner-annotated \emph{oracle-only}: independent cross-checks that
\code{compute\_cumulants} never calls but the production code is tested
against.

\begin{center}
\begin{longtable}{@{}>{\ttfamily\footnotesize\raggedright}p{0.31\textwidth} >{\raggedright}p{0.50\textwidth} >{\centering}p{0.05\textwidth}@{}}
\toprule
\normalfont\textbf{file} & \textbf{role} & \textbf{Ch.} \tabularnewline
\midrule
\endhead
\bottomrule
\endfoot
msrjd/integration/\newline spatial/full\_integrator.py
  & The production integrator. \code{diagram\_kinematic} and the batched
    Symanzik core (\code{\_momentum\_factor\_batch},
    \code{\_symanzik\_kernel\_batch}, \code{\_heat\_kernel\_x\_general}), the
    \code{grid}/\code{mc}/\code{bessel} back ends, and the assembly chain
    \code{diagram\_value\_x} $\to$ \code{diagram\_correlator\_x} $\to$
    \code{correlator\_2pt\_x}. ``One genuine integral evaluates every
    enumerated diagram.''
  & \ref{ch:spatial-core},\,\ref{ch:spatial-heatkernel} \tabularnewline
\addlinespace
msrjd/integration/\newline spatial/diagram\_descriptor.py
  & \code{diagram\_to\_cstack}: flattens a typed diagram into the canonical
    \code{CStackDiagram} edge list (frozen \code{CEdge} records), contracting
    each 2-point noise source into a single $C$ edge. Topology-blind by
    design --- no bubble/tadpole/sunset branch.
  & \ref{ch:spatial-core} \tabularnewline
\addlinespace
msrjd/integration/\newline spatial/momentum\_routing.py
  & \code{route\_momenta}: solves momentum conservation at each spatial
    vertex with SymPy's \code{linsolve}, returning the per-edge routing
    coefficients $(a_e,b_e)$ via \code{edge\_coeffs()}. The only SymPy in
    the spatial subsystem.
  & \ref{ch:spatial-core} \tabularnewline
\addlinespace
msrjd/integration/\newline spatial/spatial\_reduce.py
  & The textbook single-sample Symanzik reference:
    \code{symanzik\_matrices} ($\Lambda,N,Q$),
    \code{symanzik\_polynomials} ($\Usym=\det\Lambda$, $Q_{\text{eff}}$ via
    Schur complement), \code{momentum\_integral}. Exists so the batched
    production version can be checked against it.
  & \ref{ch:spatial-core} \tabularnewline
\addlinespace
msrjd/integration/\newline spatial/causal\_chambers.py
  & \code{causal\_chambers}: wraps the temporal pipeline's \code{\_CausalPoset}
    and \code{\_enumerate\_linear\_extensions} to tile the internal-time
    domain into smooth ordering chambers. Also a reference
    \code{integrate\_over\_chambers} (nested \code{scipy.quad}).
  & \ref{ch:spatial-core} \tabularnewline
\addlinespace
msrjd/integration/\newline spatial/heat\_kernel.py
  & The real-space propagator layer. \code{build\_spatial\_propagator} /
    \code{make\_g\_tx\_callables} and the analytic inverse Fourier transform
    of the momentum Gaussian to the heat kernel
    $(4\pi Bt)^{-d/2}\ee^{-x^2/4Bt}$ --- the trick that retired the
    $q$-grid.
  & \ref{ch:spatial-heatkernel} \tabularnewline
\addlinespace
msrjd/integration/\newline spatial/spatial\_correlator.py
  & The analytic $q\to x$ Fourier transforms used to land the diagram answer
    $C(q,\tau)$ in physical $(x,\tau)$ space, and the per-mode heat-kernel
    structure the bridge certifies against.
  & \ref{ch:spatial-heatkernel} \tabularnewline
\addlinespace
msrjd/integration/\newline spatial/spectral\_propagator.py
  & Builds the coupled multi-field free propagator $\ee^{-Mt}$ via the
    spectral (eigenprojector) decomposition of the reaction matrix $M$ --- the
    equal-diffusion coupled-tree path.
  & \ref{ch:spatial-coupled} \tabularnewline
\addlinespace
msrjd/integration/\newline spatial/dyson\_dressing.py
  & The Dyson--Duhamel series for unequal diffusion ($\mathcal{D}\ne D_0 I$):
    splits $\mathcal{D}=D_0 I+\hat{\mathcal{D}}$ and re-sums the residual
    $\hat{\mathcal{D}}|k|^2$ correction order by order.
  & \ref{ch:spatial-coupled} \tabularnewline
\addlinespace
msrjd/integration/\newline spatial/pipeline\_bridge.py
  & The 2{,}157-line spatial$\to$shared-pipeline bridge and dispatcher.
    ``Route through the shared pipeline, then certify'': substitutes
    $\Lap\to-q^2$, runs the same enumeration/typing/Phase-J path, certifies
    against the heat-kernel structure, then does the analytic $q\to x$
    transform. Hosts the \code{SPATIAL\_*} environment knobs and the memory
    guard.
  & \ref{ch:spatial-core},\,\ref{ch:spatial-coupled} \tabularnewline
\addlinespace
msrjd/integration/\newline spatial/generic\_evaluator.py
  & \emph{Oracle only.} An older Dyson-convolution evaluator that still
    branches on \code{is\_tadpole\_like()} --- the one surviving topology
    branch, kept as an independent cross-check. Not on the production path.
  & --- \tabularnewline
\addlinespace
msrjd/integration/\newline spatial/loop\_dyson.py
  & \emph{Oracle only.} Loop-level single-mode Dyson cross-check, reached
    only by its own test.
  & --- \tabularnewline
\addlinespace
msrjd/integration/\newline spatial/loop\_parametric.py
  & \emph{Oracle only.} A parametric (Feynman-parameter) loop cross-check.
  & --- \tabularnewline
\addlinespace
msrjd/integration/\newline spatial/temporal\_integrate.py
  & \emph{Oracle only.} A spatial-path temporal-integration cross-check.
  & --- \tabularnewline
\end{longtable}
\end{center}

\begin{note}
\textbf{The diagnostic helpers.} A few short members of this package are
inspection-only and never branched on in production:
\code{CStackDiagram.is\_tadpole\_like()} in
\file{diagram\_descriptor.py}, for example, classifies a diagram for
display but the production evaluator stays topology-blind. Whether a loop
momentum couples to the external $q$ is a property of the Symanzik $\Fsym$
\emph{downstream}, not of any descriptor flag --- which is the whole
``one integral, every diagram'' point.
\end{note}

\subsection{\texttt{msrjd/integration/} --- the two reference back ends}

Two files sit directly under \file{integration/}, outside both
\file{time\_domain/} and \file{spatial/}. Both predate the production
integrators and survive as references or stubs.

\begin{center}
\begin{longtable}{@{}>{\ttfamily\footnotesize\raggedright}p{0.31\textwidth} >{\raggedright}p{0.50\textwidth} >{\centering}p{0.05\textwidth}@{}}
\toprule
\normalfont\textbf{file} & \textbf{role} & \textbf{Ch.} \tabularnewline
\midrule
\endhead
\bottomrule
\endfoot
msrjd/integration/symbolic.py
  & A frequency-domain symbolic diagram-integral constructor for
    \emph{stationary} systems (the Helias \& Dahmen Ch.~9 method). A
    reference implementation, distinct from the time-domain production path.
  & \ref{ch:phasej-temporal} \tabularnewline
\addlinespace
msrjd/integration/numerical.py
  & \textbf{Stub only.} Docstring marked ``Status: Not started'' (an
    intended adaptive-quadrature / Monte-Carlo fallback). Imported by
    nothing; the \code{numerical} grep hits elsewhere are the English word.
  & --- \tabularnewline
\end{longtable}
\end{center}

% =====================================================================
\section{The driver: \texttt{pipeline/}}
% =====================================================================

\file{pipeline/} sequences the engine and exposes the authoring surface.
It splits into the master orchestrator (\file{compute.py}), the
underscore-prefixed phase modules it calls (\file{\_propagator.py},
\file{\_mean\_field*.py}, \file{\_diagrams.py}, the two caches), the
authoring API (\file{theory*.py}), the colored-noise rewriter, the result
accessors and serializers, and the \file{ui/} sub-package.

\subsection{Orchestration and the phase modules}

\begin{center}
\begin{longtable}{@{}>{\ttfamily\footnotesize\raggedright}p{0.31\textwidth} >{\raggedright}p{0.50\textwidth} >{\centering}p{0.05\textwidth}@{}}
\toprule
\normalfont\textbf{file} & \textbf{role} & \textbf{Ch.} \tabularnewline
\midrule
\endhead
\bottomrule
\endfoot
pipeline/compute.py
  & The spine. \code{compute\_cumulants(model, k, max\_ell, ...)} sequences
    the seven phases (expand $\to$ propagator $\to$ mean field $\to$ poles
    $\to$ enumerate $\to$ classify $\to$ Phase~J), branches spatial-vs-temporal,
    assembles the per-loop-order decomposition, and records wall-time. Does
    little arithmetic itself --- it is glue.
  & \ref{ch:compute-orchestration} \tabularnewline
\addlinespace
pipeline/\_propagator.py
  & The live propagator engine. Builds $K_{\text{ker}}\to K_{\text{ft}}\to
    G_{\text{ft}}$, finds the retarded poles, and computes the residue
    matrices \code{C\_mats} and the instantaneous matrix \code{D\_delta}.
  & \ref{ch:propagator} \tabularnewline
\addlinespace
pipeline/\_mean\_field.py
  & The legacy iteration saddle solver \code{solve\_mean\_field}: for
    theories whose steady state is a self-consistency closure
    $n^*_i=\phi(v^*_i)$. Iterates the rate-like variables, evaluates the
    rest on the fly.
  & \ref{ch:mean-field} \tabularnewline
\addlinespace
pipeline/\_mean\_field\_dae.py
  & The DAE multi-root solver \code{solve\_mean\_field\_dae}: sets
    $\partial_t\to0$ on the declared \code{.equation(...)} system,
    multi-start Newton for all roots, dedup, optional linear-stability
    classification, selects one root by index.
  & \ref{ch:mean-field} \tabularnewline
\addlinespace
pipeline/\_diagrams.py
  & The enumeration cache wrapper. \code{enumerate\_unique\_diagrams} wires
    the four-stage diagram pipeline (prediagrams $\to$ typed $\to$ causal $\to$
    unique) together with disk caching keyed by structure.
  & \ref{ch:type-assignment},\,\ref{ch:caching} \tabularnewline
\addlinespace
pipeline/\_expand\_cache.py
  & Disk cache for \code{FieldTheory.expand()} results, keyed by theory
    structure and Taylor order, written under
    \file{saved\_theories/<slug>/}. Avoids repeating the (up to
    $\sim$90-minute) Taylor step.
  & \ref{ch:caching} \tabularnewline
\addlinespace
pipeline/\_precompute.py
  & \code{precompute(model)}: the one-time structural pass
    (expand@order-2 $+$ propagator $+$ mean-field sanity check) that warms the
    caches and is invoked by the UI's \emph{Pre-compute} button.
  & \ref{ch:caching} \tabularnewline
\addlinespace
pipeline/\_grouped\_phase\_j.py
  & The per-prediagram grouping wrapper \code{compute\_correction\_td\_grouped}
    --- the dispatch site that hands one prediagram's typed-diagram family to
    the grouped evaluator.
  & \ref{ch:grouped-phasej} \tabularnewline
\end{longtable}
\end{center}

\subsection{Authoring: builders, compiler, templates, serialization}

\begin{center}
\begin{longtable}{@{}>{\ttfamily\footnotesize\raggedright}p{0.31\textwidth} >{\raggedright}p{0.50\textwidth} >{\centering}p{0.05\textwidth}@{}}
\toprule
\normalfont\textbf{file} & \textbf{role} & \textbf{Ch.} \tabularnewline
\midrule
\endhead
\bottomrule
\endfoot
pipeline/theory.py
  & The fluent builder classes \code{TemporalTheoryBuilder},
    \code{SpatialTheoryBuilder}, and the back-compatible \code{TheoryBuilder},
    plus \code{build()}. Accumulates \code{.physical\_field(...)},
    \code{.parameter(...)}, \code{.set\_action\_text(...)},
    \code{.equation(...)} declarations into the model dict.
  & \ref{ch:theory-spec} \tabularnewline
\addlinespace
pipeline/theory\_compiler.py
  & Text-expression $\to$ Sage-lambda compilation. Turns the human-typed
    action / mean-field / kernel / noise strings into the \code{SR}-returning
    lambdas the \code{FieldTheory} expander expects; hosts \code{\_FieldScalar}
    and the \code{ns} namespace.
  & \ref{ch:theory-spec} \tabularnewline
\addlinespace
pipeline/theory\_templates.py
  & Pre-canned action / kernel / noise templates that emit the Sage lambdas
    directly, so a user can pick a template instead of typing the action.
  & \ref{ch:theory-spec} \tabularnewline
\addlinespace
pipeline/theory\_serialize.py
  & The \file{.theory.py} \emph{text} format bridge.
    \code{render\_theory\_file} / \code{save\_theory\_to\_file} emit source
    text from a spec dict; \code{load\_spec\_from\_file} parses it back via
    the \code{ast} module (not by executing it). The round-trip behind the
    UI's Save/Load.
  & \ref{ch:theory-files} \tabularnewline
\addlinespace
pipeline/colored\_to\_markovian.py
  & The builder-level preprocessor that rewrites a colored-noise cumulant
    term into white noise by adding one auxiliary OU field per noisy field,
    so the analytic integrators apply at $\ell\ge1$.
  & \ref{ch:colored-markovian} \tabularnewline
\addlinespace
pipeline/spatial\_operator\_ir.py
  & The \code{Lap}/\code{Dt}/\code{Dx} operator intermediate representation:
    hosts spatial-derivative vertices (Model~B, Burgers, KPZ) as inert
    symbolic nodes, gives the operator algebra, and lowers them to ring
    generators so the Taylor machinery works unchanged.
  & \ref{ch:theory-spec},\,\ref{ch:spatial-coupled} \tabularnewline
\end{longtable}
\end{center}

\subsection{Result access, serialization, and reporting}

\begin{center}
\begin{longtable}{@{}>{\ttfamily\footnotesize\raggedright}p{0.31\textwidth} >{\raggedright}p{0.50\textwidth} >{\centering}p{0.05\textwidth}@{}}
\toprule
\normalfont\textbf{file} & \textbf{role} & \textbf{Ch.} \tabularnewline
\midrule
\endhead
\bottomrule
\endfoot
pipeline/access.py
  & Natural-name accessors for \code{compute\_cumulants} results: hides the
    internal \msrjd{} naming (\code{dn} for fluctuations, \code{nstar} for
    saddle values) behind the model's declared user-facing names, via
    \code{model['naming\_convention']}.
  & \ref{ch:cumulants-moments} \tabularnewline
\addlinespace
pipeline/save.py
  & NPZ $+$ CSV serialization of results with a $k$-/$\ell$-/model-adaptive
    schema: per-loop-order curves (\code{C\_tree}, \code{C\_<n>\_loop},
    \code{C\_total}) plus mean-field arrays.
  & \ref{ch:cumulants-moments} \tabularnewline
\addlinespace
pipeline/report.py
  & A multi-page PDF showing prediagrams, typed-diagram assignments, and
    per-diagram numerical contributions --- an intuition-building prototype.
  & \ref{ch:cumulants-moments} \tabularnewline
\addlinespace
pipeline/\_\_init\_\_.py
  & The package's public API surface: re-exports \code{compute\_cumulants}
    and the high-level helpers so \code{from pipeline import ...} works in
    scripts and notebooks.
  & \ref{ch:compute-orchestration} \tabularnewline
\end{longtable}
\end{center}

\subsection{\texttt{pipeline/ui/} --- the graphical theory builder}

\begin{center}
\begin{longtable}{@{}>{\ttfamily\footnotesize\raggedright}p{0.31\textwidth} >{\raggedright}p{0.50\textwidth} >{\centering}p{0.05\textwidth}@{}}
\toprule
\normalfont\textbf{file} & \textbf{role} & \textbf{Ch.} \tabularnewline
\midrule
\endhead
\bottomrule
\endfoot
pipeline/ui/main.py
  & The 3{,}071-line \code{TheoryUI}: a 10-tab \code{ipywidgets} editor with
    a live readiness sidebar and Save/Preview/Pre-compute/Load/Reset
    buttons; \code{\_collect()} snapshots the form into a spec dict for
    \code{theory\_serialize}.
  & \ref{ch:ui} \tabularnewline
\addlinespace
pipeline/ui/widgets.py
  & The reusable \code{ipywidgets} primitives the tabs compose from:
    \code{DynamicTable}, \code{vector\_input}, \code{matrix\_input},
    \code{expression\_input}, \code{textarea\_input}, \code{paste\_button},
    \code{read\_system\_clipboard}.
  & \ref{ch:ui} \tabularnewline
\addlinespace
pipeline/ui/\_\_init\_\_.py
  & Re-exports \code{TheoryUI}.
  & \ref{ch:ui} \tabularnewline
\end{longtable}
\end{center}

\subsection{\texttt{pipeline/theories/} and \texttt{pipeline/examples/}}

These are \emph{hand-written} model dicts and example scripts that predate
the \file{theories/*.theory.py} format. They are not superseded --- the
demo notebook and several tests still import them directly.

\begin{center}
\begin{longtable}{@{}>{\ttfamily\footnotesize\raggedright}p{0.31\textwidth} >{\raggedright}p{0.50\textwidth} >{\centering}p{0.05\textwidth}@{}}
\toprule
\normalfont\textbf{file} & \textbf{role} & \textbf{Ch.} \tabularnewline
\midrule
\endhead
\bottomrule
\endfoot
pipeline/theories/\newline linear\_hawkes\_2pop\_delta.py
  & Hand-written model dict: linear Hawkes, 2 populations, instantaneous
    ($\delta$) synaptic kernel. Equivalent to
    \file{models/hawkes\_linear\_sage.py}.
  & \ref{ch:theory-spec} \tabularnewline
\addlinespace
pipeline/theories/\newline linear\_hawkes\_2pop\_expg.py,
  \newline linear\_hawkes\_2pop\_expg\_gtas.py,
  \newline quad\_hawkes\_2pop\_expg.py,
  \newline quad\_hawkes\_2pop\_expg\_gtas.py
  & The exponential-kernel and GTaS (Gaussian-truncated-and-shifted)
    Hawkes variants, linear and quadratic. Imported by
    \file{pipeline\_demo.ipynb}, several \file{*\_sim\_compare} notebooks, and
    \file{pipeline/tests/test\_theory\_equivalence.py}.
  & \ref{ch:theory-spec} \tabularnewline
\addlinespace
pipeline/examples/\newline run\_linear\_gtas.py
  & A standalone end-to-end example script: compute the $k=2$
    cross-correlator $\langle\delta n_1(0)\,\delta n_2(\tau)\rangle$ for the
    linear-GTaS model and emit a PDF report.
  & \ref{ch:compute-orchestration} \tabularnewline
\addlinespace
pipeline/tests/\newline test\_theory\_equivalence.py
  & Asserts the hand-written \file{pipeline/theories/} dicts and the
    builder-produced theories agree.
  & --- \tabularnewline
\end{longtable}
\end{center}

% =====================================================================
\section{The validators: \texttt{models/}}
% =====================================================================

\file{models/} is the trust chain's other half: completely independent
direct simulators and the estimators that turn their time series into the
same cumulant slices the engine predicts. The two are overlaid in the
\file{*\_sim\_compare} notebooks; if the diagrammatics are correct, the
curves coincide within Monte-Carlo error.

\subsection{SDE and field simulators}

\begin{center}
\begin{longtable}{@{}>{\ttfamily\footnotesize\raggedright}p{0.31\textwidth} >{\raggedright}p{0.50\textwidth} >{\centering}p{0.05\textwidth}@{}}
\toprule
\normalfont\textbf{file} & \textbf{role} & \textbf{Ch.} \tabularnewline
\midrule
\endhead
\bottomrule
\endfoot
models/ou\_langevin\_sim\_numba.py
  & Euler--Maruyama simulators for scalar Langevin SDEs (OU-type theories),
    \code{numba}-JIT compiled. The simulator behind the quickstart's
    \code{sim\_ou\_quartic\_numba}.
  & \ref{ch:simulators} \tabularnewline
\addlinespace
models/spatial\_field\_1d\_sim.py
  & Spectral exponential-Euler (ETD1) simulator for a 1D scalar stochastic
    field on a periodic ring --- the spatial-v1 cross-check.
  & \ref{ch:simulators} \tabularnewline
\addlinespace
models/spatial\_field\_2d\_sim.py
  & 2D spectral ETD1 field simulator on a periodic square --- the $d=2$
    validation oracle.
  & \ref{ch:simulators} \tabularnewline
\addlinespace
models/spatial\_field\_3d\_sim.py
  & 3D spectral ETD1 field simulator on a periodic cube --- the $d=3$
    validation oracle.
  & \ref{ch:simulators} \tabularnewline
\addlinespace
models/spatial\_field\_phi6\_1d\_sim.py
  & ETD1 simulator for the 1D stochastic $\phi^6$ field (Allen--Cahn $+$
    quintic force) --- the cross-check for the quintic theory.
  & \ref{ch:simulators} \tabularnewline
\addlinespace
models/coupled\_rd\_1d\_sim.py
  & $N$-species coupled stochastic reaction--diffusion Langevin simulator on
    a ring --- the coupled-field cross-check.
  & \ref{ch:simulators} \tabularnewline
\addlinespace
models/dendritic\_linear\_sim\_numba.py
  & Soma$+$dendrite compartmental linear-rate simulator (\code{numba} JIT)
    for the dendritic theory.
  & \ref{ch:simulators} \tabularnewline
\end{longtable}
\end{center}

\subsection{Hawkes / point-process simulators}

\begin{center}
\begin{longtable}{@{}>{\ttfamily\footnotesize\raggedright}p{0.31\textwidth} >{\raggedright}p{0.50\textwidth} >{\centering}p{0.05\textwidth}@{}}
\toprule
\normalfont\textbf{file} & \textbf{role} & \textbf{Ch.} \tabularnewline
\midrule
\endhead
\bottomrule
\endfoot
models/hawkes\_sim\_numba.py
  & \code{numba}-JIT Euler-step simulator for the linear Hawkes process
    (fixed-step Poisson-draw realization).
  & \ref{ch:simulators} \tabularnewline
\addlinespace
models/hawkes\_sim\_multipop\_numba.py
  & The linear \emph{multipopulation} Hawkes stepper.
  & \ref{ch:simulators} \tabularnewline
\addlinespace
models/hawkes\_sim\_expg\_numba.py,
  \newline hawkes\_sim\_quad\_expg\_numba.py
  & Exponential-kernel Hawkes simulators, linear and quadratic.
  & \ref{ch:simulators} \tabularnewline
\addlinespace
models/\newline hawkes\_sim\_linear\_expg\_gtas\_numba.py,
  \newline hawkes\_sim\_quad\_expg\_gtas\_numba.py
  & The GTaS-kernel Hawkes simulators, linear and quadratic.
  & \ref{ch:simulators} \tabularnewline
\end{longtable}
\end{center}

\subsection{Estimators and hand-written Hawkes model dicts}

\begin{center}
\begin{longtable}{@{}>{\ttfamily\footnotesize\raggedright}p{0.31\textwidth} >{\raggedright}p{0.50\textwidth} >{\centering}p{0.05\textwidth}@{}}
\toprule
\normalfont\textbf{file} & \textbf{role} & \textbf{Ch.} \tabularnewline
\midrule
\endhead
\bottomrule
\endfoot
models/cumulant\_estimator.py
  & The general $k$-point connected-cumulant estimator
    (\code{estimate\_kpoint\_slices}): bins a trajectory and estimates the
    connected cumulant density versus lag, removing the discrete-spike
    shot-noise diagonal.
  & \ref{ch:simulators} \tabularnewline
\addlinespace
models/cumulant\_direct\_numba.py
  & \code{numba}-JIT direct (non-FFT) 3-point cumulant estimator --- the
    $k=3$ companion to the general estimator.
  & \ref{ch:simulators} \tabularnewline
\addlinespace
models/hawkes\_sage.py,
  \newline hawkes\_linear\_sage.py
  & Hand-written SageMath model specifications for the nonlinear and linear
    Hawkes processes.
  & \ref{ch:theory-spec} \tabularnewline
\addlinespace
models/hawkes\_linear\_expg.py,
  \newline hawkes\_quad\_expg.py,
  \newline hawkes\_linear\_expg\_gtas.py,
  \newline hawkes\_quad\_expg\_gtas.py
  & Hand-written model dicts for the exponential- and GTaS-kernel Hawkes
    families (the symbolic counterparts of the simulators above).
    \file{hawkes\_quad\_expg.py} is the canonical worked example in the
    field-theory-core chapter.
  & \ref{ch:fieldtheory-core},\,\ref{ch:theory-spec} \tabularnewline
\end{longtable}
\end{center}

% =====================================================================
\section{The models: \texttt{theories/}}
% =====================================================================

Each \file{theories/<name>.theory.py} is one physical model: an importable
module exposing exactly \code{build()} (returns the model dict),
\code{DEFAULT\_FUNDAMENTAL} (default numeric parameters), and
\code{METADATA} (run recommendations). \code{dd.load\_theory('<name>')}
imports the file and calls \code{build()}. Rather than list all
forty-plus, the table groups them by the builder they use and the
capability each exercises; the loadable name is the filename without the
\file{.theory.py} suffix.

\begin{center}
\begin{longtable}{@{}>{\raggedright}p{0.30\textwidth} >{\ttfamily\footnotesize\raggedright}p{0.40\textwidth} >{\centering}p{0.16\textwidth}@{}}
\toprule
\textbf{group (builder)} & \normalfont\textbf{representative files} & \textbf{Ch.} \tabularnewline
\midrule
\endhead
\bottomrule
\endfoot
Temporal, single field, the nonlinear baseline
  & ou\_quartic, ou\_quartic\_double\_well, ou\_sextic,
    toy\_quartic\_double\_well
  & \ref{ch:quickstart},\,\ref{ch:theory-files} \tabularnewline
\addlinespace
Temporal, colored / correlated noise
  & ou\_quartic\_colored, ou\_quartic\_two\_dim,
    ou\_quartic\_two\_dim\_corr, ou\_quartic\_two\_dim\_color\_corr
  & \ref{ch:colored-markovian} \tabularnewline
\addlinespace
Temporal Hawkes (point process)
  & linear\_hawkes, quadratic\_hawkes, quadratic\_hawkes\_alpha
  & \ref{ch:theory-spec} \tabularnewline
\addlinespace
Temporal neuron / compartmental tests
  & single\_pop\_dendritic\_linear,
    single\_population\_conductance\_test,
    single\_population\_spike\_reset\_test,
    single\_population\_quad\_exp\_test, multipopulation\_test
  & \ref{ch:theory-spec},\,\ref{ch:mean-field} \tabularnewline
\addlinespace
Spatial, single field, plain vertex
  & allen\_cahn\_1d\_subcritical\_infinite,
    allen\_cahn\_1d\_subcritical\_pbc,
    allen\_cahn\_quintic\_1d\_subcritical\_infinite,
    reaction\_diffusion\_quadratic\_1d
  & \ref{ch:spatial-core},\,\ref{ch:spatial-heatkernel} \tabularnewline
\addlinespace
Spatial, derivative vertex (operator IR)
  & kpz\_1d, kpz\_2d, burgers\_1d, edwards\_wilkinson\_1d,
    reaction\_diffusion\_conserved\_1d
  & \ref{ch:spatial-heatkernel},\,\ref{ch:spatial-coupled} \tabularnewline
\addlinespace
Spatial, coupled / combined fields
  & coupled\_rd\_2species\_1d, reaction\_diffusion\_2d,
    combined\_allencahn\_modelb\_kpz\_\{1,2,3\}d, linear\_field\_2d,
    linear\_diffusion\_test
  & \ref{ch:spatial-coupled} \tabularnewline
\end{longtable}
\end{center}

\begin{note}
\textbf{Builder-class hint.} Files using \code{TemporalTheoryBuilder} are
time-only; \code{SpatialTheoryBuilder} files are PDEs; a handful (e.g.
\file{combined\_allencahn\_modelb\_kpz\_1d.theory.py},
\file{linear\_field\_2d.theory.py}) use the back-compatible
\code{TheoryBuilder} that auto-detects at \code{build()} time. The
\code{*\_test} files are regression fixtures exercising one feature each
(spike-reset resets, conductance synapses, $\delta$- vs exponential
voltage), not headline physics.
\end{note}

% =====================================================================
\section{The front desk and notebooks: \texttt{notebooks/}}
% =====================================================================

\subsection{The convenience layer}

\begin{center}
\begin{longtable}{@{}>{\ttfamily\footnotesize\raggedright}p{0.31\textwidth} >{\raggedright}p{0.50\textwidth} >{\centering}p{0.05\textwidth}@{}}
\toprule
\normalfont\textbf{file} & \textbf{role} & \textbf{Ch.} \tabularnewline
\midrule
\endhead
\bottomrule
\endfoot
notebooks/daedalus.py
  & The 1{,}101-line front desk imported as \code{dd}. The four verbs
    (\code{load\_theory}, \code{Config}, \code{run}, \code{plot\_cumulant})
    plus introspection helpers and the \code{output='moment'} assembly
    layer. Thin orchestration --- it does almost no physics.
  & \ref{ch:quickstart},\,\ref{ch:engine-api},\,\ref{ch:cumulants-moments} \tabularnewline
\end{longtable}
\end{center}

\subsection{The notebook directories}

The notebooks are not modules to read line-by-line; they are the
runnable surface. They are grouped by purpose:

\begin{center}
\begin{longtable}{@{}>{\ttfamily\footnotesize\raggedright}p{0.31\textwidth} >{\raggedright}p{0.50\textwidth} >{\centering}p{0.05\textwidth}@{}}
\toprule
\normalfont\textbf{directory} & \textbf{role} & \textbf{Ch.} \tabularnewline
\midrule
\endhead
\bottomrule
\endfoot
notebooks/examples/
  & The nine bonafide example notebooks, one per pipeline capability,
    each with a model-info $\to$ theory $\to$ sim structure
    (\file{temporal\_ou\_quartic\_white.ipynb} is the quickstart's).
  & \ref{ch:quickstart},\,\ref{ch:engine-api} \tabularnewline
\addlinespace
notebooks/templates/
  & The four group templates (\{temporal,spatial\}$\times$\{single,multi\})
    plus \code{*\_sim\_compare} variants that every demo is generated from.
  & \ref{ch:engine-api} \tabularnewline
\addlinespace
notebooks/temporal/
  & The temporal \code{*\_sim\_compare} validation notebooks (OU quartic,
    Hawkes, conductance, spike-reset, $\dots$) --- theory overlaid on
    simulation.
  & \ref{ch:simulators} \tabularnewline
\addlinespace
notebooks/spatial/
  & The spatial \code{*\_sim\_compare} validation notebooks (Allen--Cahn,
    KPZ, Burgers, Model~B, coupled RD, $d=1,2,3$).
  & \ref{ch:simulators} \tabularnewline
\addlinespace
notebooks/legacy/
  & The pre-pipeline demo notebooks (enumeration, field-theory-Sage,
    Hawkes time-domain), kept for provenance.
  & --- \tabularnewline
\addlinespace
notebooks/\newline theory\_builder.ipynb,
  \newline theory\_runner.ipynb
  & The single-cell launchers for the \code{TheoryUI} form and for running a
    saved \file{.theory.py}.
  & \ref{ch:ui} \tabularnewline
\end{longtable}
\end{center}

% =====================================================================
\section{Supporting trees: tests, scripts, legacy, docs}
% =====================================================================

The remaining top-level directories are not part of the runtime data flow
but are essential to working in the repository.

\begin{center}
\begin{longtable}{@{}>{\ttfamily\footnotesize\raggedright}p{0.27\textwidth} >{\raggedright}p{0.54\textwidth} >{\centering}p{0.05\textwidth}@{}}
\toprule
\normalfont\textbf{directory} & \textbf{role} & \textbf{Ch.} \tabularnewline
\midrule
\endhead
\bottomrule
\endfoot
tests/
  & The pytest suite (61 \code{test\_*.py} files). One file per subsystem
    plus the oracle cross-checks: e.g. \file{test\_full\_integrator.py},
    \file{test\_spatial\_reduce.py}, \file{test\_momentum\_routing.py},
    \file{test\_symmetry.py}, \file{test\_grouped\_vs\_perdiag.py},
    \file{test\_all\_k\_boltzmann.py}, \file{test\_kpz\_burgers\_sim.py},
    \file{test\_generic\_evaluator.py} (the spatial oracle),
    \file{test\_fork\_safety.py}.
  & all \tabularnewline
\addlinespace
scripts/
  & Standalone runnable scripts; currently
    \file{scripts/profile\_pipeline\_k3.py}, a $k=3$ profiling harness.
  & --- \tabularnewline
\addlinespace
legacy/
  & The pre-pipeline \code{sympy}-era codebase (the original Enumeration
    Code, Field Theory Framework, Theory Builder, and \code{*\_sympy.py}
    modules), isolated and kept for provenance. \emph{Not} on any live
    path.
  & --- \tabularnewline
\addlinespace
saved\_theories/
  & The on-disk \code{.sobj} cache (per-theory \file{<slug>/} directories of
    pickled propagators, expanded actions, and unique-diagram sets, plus a
    \code{manifest.json}). Gitignored, regenerable.
  & \ref{ch:caching} \tabularnewline
\addlinespace
docs/
  & The prose documentation, design records, and this manual
    (\file{docs/manual/}). \file{docs/spatial\_pipeline.md} is the
    authoritative spatial-state reference; \file{docs/CHANGELOG.md} is the
    chronological fix log.
  & --- \tabularnewline
\addlinespace
scratch/, \newline pipeline\_outputs/
  & Throwaway probe scripts/logs and generated run artefacts respectively.
    Not part of the codebase; safe to delete. (\file{scratch/} also holds
    the tropical-Monte-Carlo research PDFs cited by the efficiency roadmap.)
  & --- \tabularnewline
\end{longtable}
\end{center}

% =====================================================================
\section{The shortest path to the file you want}
% =====================================================================

To close, a reverse index --- start from what you are trying to do, and
this tells you which file (and chapter) to open:

\begin{description}[style=nextline,leftmargin=1.4em]
  \item[``I want to add a new model.''] Write a
        \file{theories/<name>.theory.py} using \file{pipeline/theory.py}'s
        builder (Chapter~\ref{ch:theory-spec},
        Chapter~\ref{ch:theory-files}), or point-and-click it with
        \file{pipeline/ui/main.py} (Chapter~\ref{ch:ui}).
  \item[``I want to run a model and plot it.''] \file{notebooks/daedalus.py}
        --- the \code{load\_theory $\to$ Config $\to$ run $\to$ plot} spine
        (Chapter~\ref{ch:quickstart}, Chapter~\ref{ch:engine-api}).
  \item[``The answer is wrong and I suspect the diagrams.''] The enumeration
        is \file{msrjd/enumeration/loop\_diagram\_enumeration.py}
        (Chapter~\ref{ch:enumeration}); the typing/symmetry is
        \file{msrjd/diagrams/} (Chapter~\ref{ch:type-assignment}).
  \item[``The answer is wrong and I suspect the loop integral.''] Temporal:
        \file{msrjd/integration/time\_domain/final\_integral.py}
        (Chapter~\ref{ch:phasej-temporal}). Spatial:
        \file{msrjd/integration/spatial/full\_integrator.py}
        (Chapter~\ref{ch:spatial-core},
        Chapter~\ref{ch:spatial-heatkernel}).
  \item[``The saddle looks wrong.''] \file{pipeline/\_mean\_field\_dae.py}
        (Chapter~\ref{ch:mean-field}).
  \item[``It is too slow / it recomputes every run.''] The caches:
        \file{pipeline/\_expand\_cache.py} and \file{pipeline/\_diagrams.py}
        (Chapter~\ref{ch:caching}).
  \item[``I want to check it against a simulation.''] The matching
        \file{models/*\_sim*.py} simulator $+$
        \file{models/cumulant\_estimator.py} (Chapter~\ref{ch:simulators}),
        driven from a \file{notebooks/\{temporal,spatial\}/*\_sim\_compare.ipynb}.
  \item[``Where does it all start?''] \code{compute\_cumulants} in
        \file{pipeline/compute.py} (Chapter~\ref{ch:compute-orchestration})
        --- the one function every path runs through.
\end{description}
